\chapter{Interpretabilidad}
\label{ch:interpretabilidad}

Este capítulo corresponde al notebook \texttt{07\_interpretabilidad} y al
script \texttt{run\_interpretability.py}. Se emplean
LIME~\cite{ribeiro2016lime} y SHAP~\cite{lundberg2017shap} para explicar
predicciones de mBERT.

\section{Falsos positivos y jerga}

La Tabla~\ref{tab:fp-jerga} cuenta falsos positivos de la clase No Tóxico
según presencia de jerga. Hay 12 FP con jerga y 15 sin jerga; dado el menor
tamaño del grupo con jerga, la tasa relativa es mayor (coherente con la sonda
del Capítulo~\ref{ch:evaluacion}).

\begin{table}[H]
\centering
\caption{Falsos positivos No Tóxico según jerga.}
\label{tab:fp-jerga}
\begin{tabular}{lc}
\toprule
Grupo & FP \\
\midrule
fp\_con\_jerga & 12 \\
fp\_sin\_jerga & 15 \\
\bottomrule
\end{tabular}

\end{table}

\section{Concordancia LIME--SHAP}

Se calcula el índice de Jaccard entre los top-$k$ tokens de LIME y SHAP por
instancia (Tabla~\ref{tab:jaccard}). En la corrida actual el Jaccard es
\textbf{0} en todas las instancias muestreadas. Esto puede deberse a
desalineación de tokenización, fallos parciales de SHAP o presupuestos de
evaluación bajos (\texttt{max\_evals}). Se reporta como limitación: las
explicaciones LIME siguen siendo informativas de forma aislada, pero no se
valida solapamiento con SHAP en esta ejecución.

\begin{table}[H]
\centering
\caption{Jaccard top-$k$ entre LIME y SHAP por instancia.}
\label{tab:jaccard}
\begin{tabular}{l}
\toprule
Nota \\
\midrule
Jaccard LIME--SHAP omitido: solapamiento nulo tras alineación WordPiece (ver nota en capítulo) \\
\bottomrule
\end{tabular}

\end{table}

\section{Importancias globales LIME}

Cuando el agregado de pesos LIME está disponible, la Tabla~\ref{tab:lime} y la
Figura~\ref{fig:lime} muestran los tokens más influyentes por clase. Si el
archivo de métricas está vacío (corrida incompleta), las tablas/figuras
incluyen un marcador de datos no disponibles.

\begin{table}[H]
\centering
\caption{Top tokens LIME por clase (agregado).}
\label{tab:lime}
\begin{tabular}{l}
\toprule
Nota \\
\midrule
LIME global weights no disponibles en esta corrida \\
\bottomrule
\end{tabular}

\end{table}

\begin{figure}[H]
\centering
\includegraphics[width=0.9\textwidth]{08_lime_global_weights.png}
\caption{Importancias globales LIME por clase.}
\label{fig:lime}
\end{figure}
